\chapter{METODOLOGI PENELITIAN}
\label{bab:metodologi}

\section{Desain Penelitian}
\label{sec:desain}

Dummy konten desain penelitian. Bagian ini menjelaskan pendekatan dan desain penelitian yang digunakan.

\section{Data dan Sumber Data}
\label{sec:data}

Dummy konten data dan sumber data. Bagian ini menjelaskan karakteristik dataset yang digunakan.

\subsection{Dataset A}
\label{subsec:dataset_a}

Deskripsi Dataset A dengan karakteristik sebagai berikut:

\begin{itemize}
    \item Jumlah sampel: 1000
    \item Jumlah kelas: 7
    \item Sumber: sumber data
\end{itemize}

\subsection{Dataset B}
\label{subsec:dataset_b}

Deskripsi Dataset B dengan karakteristik sebagai berikut:

\begin{itemize}
    \item Jumlah sampel: 2000
    \item Jumlah kelas: 7
    \item Sumber: sumber data
\end{itemize}

\section{Variabel Penelitian}
\label{sec:variabel}

Dummy konten variabel penelitian.

\begin{table}[h]
    \centering
    \caption{Variabel Penelitian}
    \label{tab:variabel}
    \begin{tabularx}{\textwidth}{X X}
        \toprule
        \textbf{Variabel Bebas} & \textbf{Variabel Terikat} \\
        \midrule
        Variabel X1 & Variabel Y \\
        Variabel X2 & \\
        \bottomrule
    \end{tabularx}
\end{table}

\section{Alat dan Bahan}
\label{sec:alat_bahan}

\subsection{Spesifikasi Perangkat}
\label{subsec:perangkat}

\begin{itemize}
    \item Prosesor: Intel Core i7
    \item RAM: 32 GB
    \item GPU: NVIDIA RTX 3080
    \item Storage: SSD 1 TB
\end{itemize}

\subsection{Perangkat Lunak}
\label{subsec:perangkat_lunak}

\begin{itemize}
    \item Python 3.10
    \item PyTorch 2.0
    \item CUDA 12.1
\end{itemize}

\section{Prosedur Penelitian}
\label{sec:prosedur}

Dummy konten prosedur penelitian. Gambar \ref{fig:prosedur} menunjukkan diagram alur prosedur penelitian.

\begin{figure}[h]
    \centering
    \fbox{\parbox[c][6cm][c]{12cm}{\centering Diagram Alur Prosedur Penelitian}}
    \caption{Diagram Alur Prosedur Penelitian}
    \label{fig:prosedur}
\end{figure}

\section{Teknik Pengumpulan Data}
\label{sec:teknik_pengumpulan}

Dummy konten teknik pengumpulan data.

\section{Teknik Analisis Data}
\label{sec:teknik_analisis}

Dummy konten teknik analisis data.

\subsection{Metrik Evaluasi}
\label{subsec:metrik}

Metrik yang digunakan untuk evaluasi model:

\begin{equation}
    Accuracy = \frac{TP + TN}{TP + TN + FP + FN}
    \label{eq:accuracy}
\end{equation}

\begin{equation}
    F1 = \frac{2 \times Precision \times Recall}{Precision + Recall}
    \label{eq:f1}
\end{equation}